\begin{center}
\LARGE
Towards modelling the topology of homogeneous manifolds
              by means of symbolic computation 

\end{center}

\begin{center}
\Large
Michael Joswig

\end{center}

\noindent
Date:  July 17th (Wednesday)\\
Time:  11:10-11:35\\

\begin{center}
\large
Abstract
\end{center}

The theory of Lie groups is among the central topics in modern
mathematics. It represents a natural link between geometry, group
theory and mathematical physics. Problems arising from physics,
e.g.~from the theory of special relativity, are transferable to the
language of geometry.  Often the geometrical questions can be answered
by group-theoretical means in case the geometrical structure has
enough symmetry, i.e. it admits sufficiently many automorphisms. In
this respect, transitive actions are most important. Manifolds which
admit a transitive action of some Lie group are called {\it homogeneous}.

In recent years quite a few attempts have been made to bring Lie
theory to the computer, e.g. see van Leeuwen, Cohen and Lisser [.Cohen
LiE.], Moody, Patera and Rand [.simpLie.], B�di and Joswig [.boedi
Joswig.], not to mention the various packages devoted to the study of
finite groups of Lie type. When it comes to solving problems arising
from real Lie groups by means of a computer frequently the derived
linearized problems are studied instead. This means one tries to
translate the initial problem about Lie groups into a (usually simpler
one) about Lie algebras.  However, the correspondence between Lie
groups and Lie algebras is {\it not} 1--1. In fact, for a given Lie algebra
g there are usually several (not isomorphic) groups having g as their
associated Lie algebra. They are only {\it locally} isomorphic, which means
that they are algebraically somehow similar but differ with respect to
their global topological properties. Roughly speaking, during the
linearization process the global topological information is lost.

Here we want to suggest an approach with is in some sense
complementary. Instead of skipping the topological information we
intend to focus on it.

We want to look for a computational method to solve problems of the
following kind and related ones: Given a Lie group $\Gamma$ and a
homogeneous manifold $M$. Does there exist a (continuous) transitive
action of $\Gamma$ on $M$?

Without any doubt this immediately raises at least two questions. How
should a non-discrete Lie group (being of uncountable cardinality as a
set) be represented in the computer? How can one even expect that a
question like this is decidable at all? Instead of trying to solve the
problem in this generality we want to discuss a certain approximation
of the problem which can be outlined as follows. A transitive action
of $\Gamma$ on $M$ gives rise to an infinitely long exact sequence of their
respective homotopy groups (see below). The group $\Gamma$ and the
manifold $M$ are described by some of their algebraic and topological
properties such as the dimension, information about compactness, maybe
some homotopy groups which are known etc. It is explicitly admitted
that only partial information will be given. Assume that a transitive
action exists. Select finitely many among the infinitely many pieces
of information obtained from the long exact homotopy sequence.
Translate this into an arithmetical problem. Check for a
contradiction. In case we actually arrive at a contradiction then this
falsifies the assumption concerning the existence of a transitive
action.  However, if we do not obtain a contradiction, we might not
gain any information.

